\section{Evaluation}
\label{sec:storyevaluation}
All experimental results remain unfilled. The evaluation distinguishes service
fairness, conditional advantage over CFL, incremental policy cost, and
application benefit.

\draftslot{EVAL-SETUP}{Specify revisions, machine/compiler, placement and total
core budget, workload provenance, repetitions, uncertainty, and instrumentation
cost.}

\subsection{Does the scheduler distribute the intended service?}
Heterogeneous operations test whether balanced accounting corresponds to
balanced useful service. Traces relate selectable membership and every
charge to completed work, comparing billed and useful-service spread
(Section~\ref{sec:storyfairness}). The initial test
uses equal credits and a fixed continuously selectable cohort with bounded
nonnegative increments and no resets. Separate arrival/reactivation tests
examine admission rather than borrowing that guarantee. Agreement cannot
establish unconditional correctness of current FC-PQ.

\draftslot{EVAL-SERVICE}{Supply cohort and accounting evidence, per-class
service, spread trajectories, and admission failures; identify which contract
premise each violation challenges.}

\subsection{When does FC-PQ outperform CFL?}
The principal comparison will use CFL's specified fair-policy mode, aligning
the scheduled entity and service metric with FC-PQ. Offered work, completed
operation mix, and hardware resources must be matched; per-class completions
and useful service prevent cheap operations from inflating throughput.
Varying footprint, load, and placement tests
Equation~\eqref{eq:storycrossover}, including CFL's NUMA locality.
Actual completions, not queue pops or budgets, determine batch occupancy.

We will calibrate costs independently and predict batch occupancy and executor
changes before held-out runs in the model's stated regime. Measured traces
will diagnose errors, not count as predictions (Appendix~\ref{app:costs}).

\placeholderfigure{EVAL-CFL}{Matched-work CFL and FC-PQ throughput across
footprint/load regimes, with useful-service fairness alongside. Mark predicted
crossovers, uncertainty, and the region where FC-PQ loses.}{Planned comparison
of conditional performance advantage and service allocation.}

\subsection{What is the incremental price of fairness?}
We will compare fair and policy-disabled versions of the same delegation
protocol; original FC versus FC-PQ is not this control. Separate controlled
measurements will distinguish direct accounting and selection costs from
changes in order, batching, locality, and completed mix
(Appendix~\ref{app:evidence}).

\draftslot{EVAL-COST}{Supply matched-control costs, independent calibration,
held-out prediction errors and crossover failures, distinguishing direct
components from indirect execution changes.}

\subsection{Does the application retain the benefit?}
An application comparison must preserve semantics and return to heterogeneous
service demand, counting adaptation and marshalling costs. Per-class return
distributions will separate executor-callers from waiting callers: completion
need not mean the executor can return immediately. Sparse activity, capacity
eviction, large request metadata, and tight return budgets deliberately probe
where locality or amortization ceases to help.

\draftslot{EVAL-APPLICATION}{Supply a semantically equivalent application,
delegate ownership/affinity restrictions, useful-work throughput and role-aware
return evidence, including losing regimes and non-lock bottlenecks.}
